% RESEARCH_DRAFT_SHA256: 0ab57b1c0bfb9d27adecd7922424462da4b6aafb795f2037d986a864ff5c2e51
% Layout and source-reference wrappers expanded for the web renderer.
\documentclass[11pt]{article}
\usepackage[margin=25mm]{geometry}
\usepackage{amsthm}
\usepackage{fontspec,unicode-math}
\setmainfont{Latin Modern Roman}
\setmathfont{Latin Modern Math}
\usepackage{amsmath,amssymb,mathtools}
\usepackage{xurl,hyperref}
\hypersetup{colorlinks=true,urlcolor=blue,linkcolor=blue}
\setcounter{secnumdepth}{0}
\setlength{\parindent}{0pt}
\setlength{\parskip}{5pt}
\setlength{\emergencystretch}{3em}
\newcommand{\sref}[2]{\hyperlink{source:#1:#2}{[S#2]}}
\newcommand{\cataloguescope}[1]{\paragraph{Source question.}#1}
\newcommand{\floor}[1]{\left\lfloor#1\right\rfloor}
\newtheorem{proposition}{Proposition}
\newtheorem{lemma}[proposition]{Lemma}
\begin{document}
\paragraph{Research manuscript.} OpenAI \texttt{gpt-6-astra}, reasoning \texttt{ultra}; 11 October 2026.
\paragraph{Preserved source body.} The complete catalogue mathematical text and sources follow with layout and source links expanded, including source qualifications and its historical preparation notes. New research begins at the section headed Research attempt.

\section{Sample Complexity of Vertex-Distribution-Free Graph Testing}\label{statement}

\subsection{Definitions and mathematical statement}

\paragraph{Shared notation.}Standard number systems are denoted by $\mathbb N,\mathbb Z,\mathbb Q,\mathbb R,\mathbb C$; index ranges, computational models and specialized notation follow the cited problem source and the definitions below.


A graph property is hereditary if it is preserved by taking induced subgraphs, and extendable if every graph having it can be enlarged by one vertex while retaining the property. In the vertex-distribution-free model a tester samples vertices independently from an unknown distribution $D$ and queries adjacency. Distance is measured by the $D(u)D(v)$-weighted mass of changed adjacency pairs. A tester distinguishes property members from graphs at distance at least $\varepsilon$ with fixed bounded error. The standard dense model uses the uniform vertex distribution.
\paragraph{Statement recorded in the supplied source.}
Does every extendable hereditary graph property admit a vertex-distribution-free tester whose sample complexity is asymptotically comparable to that of a tester in the standard dense-graph model? \href{https://doi.org/10.4171/owr/2019/19}{[S2]}
\paragraph{Scope qualification.}
The comparison is an asymptotic dependence on the proximity parameter, with the property fixed. The report does not prescribe a universal comparison function.
\paragraph{Source status.}
The supplied catalogue labels this record open. Its source-status metadata was checked against the supplied snapshot on 2026-10-11. This import does not independently establish current open status or mathematical correctness. The catalogue literature assessment is dated 2026-08-22.
\subsection{Short English statement}
Compare the sampling cost of testing hereditary properties under arbitrary vertex sampling and under uniform sampling.
\subsection{Sources}
\begin{itemize}
\item[S1] Ulam AI and the UnsolvedMath Contributors, UnsolvedMath: A Curated Collection of Open Mathematics Problems, catalogue entry. Catalogue curation and metadata: CC BY 4.0. The supplied source question is retained; definitions and the English summary are editorial additions.
\url{https://huggingface.co/datasets/ulamai/UnsolvedMath}.
\item[S2] Combinatorics, Probability and Computing (2020), Oberwolfach Reports, DOI 10.4171/owr/2019/19, pp. 1113--1189
\url{https://doi.org/10.4171/owr/2019/19}.
\item[S3] Combinatorics, Probability and Computing, Oberwolfach Reports 19/2019, primary report PDF
\url{https://ems.press/content/serial-article-files/46798}.
\textit{Underlying Oberwolfach Reports material retains its source-specific licence, including CC BY-SA 4.0 where applicable; see THIRD\_PARTY\_NOTICES.md.}
\end{itemize}
\paragraph{Preparation and review.}
Original question attributed to the cited source and UnsolvedMath contributors. Codex (GPT-6; exact serving revision not exposed) assisted with definitions, English summaries, source inspection and formatting. No new human mathematical authorship or independent proof verification is asserted. Catalogue identity checked against the complete supplied registry and existing repository statements; known flagged repetitions and obvious identical conjectures excluded. Full mathematical-equivalence/literature review has not been performed. Statement preparation and final review followed contribution prompts C and D; automated checks verify structure and identity, not open status.
\label{end:statement}

\section{Research attempt}
\subsection{Literature, source context and scope}
The full original catalogue body and source qualifications are preserved above. The relevant primary workshop passages were read in the cached publisher report. Historical claims and subsequent attributed results are distinguished from the mathematical work below. A restricted example, standard lemma or known theorem is not counted as new progress on the entire original target. Missing assumptions and repeated or rhetorical extracts are not silently turned into different solved problems.

\subsection{An exact easy property and a failed general transfer}
The source defines distance using unordered changed pairs, with weight $D(u)D(v)$, and gives a full qualitative characterization before asking about sample complexity. We retain that normalization. Consider the extendable hereditary property of having no edges. Put
\[
w=\sum_{\{u,v\}\in E(G)}D(u)D(v).
\]
This is exactly the distance to the property, since every edge must be deleted. Draw $t$ independent pairs of vertices from $D$, query whether each pair is an edge, and reject upon seeing an edge.
\begin{proposition}
This tester has one-sided error and rejects an $\varepsilon$-far graph with probability at least $2/3$ when $t\ge\lceil\log(3)/(2\varepsilon)\rceil$.
\end{proposition}
\begin{proof}
A sampled ordered pair is an edge with probability $2w$: the two orders of each unordered edge have equal weight, and loops never witness an edge. Independent pairs therefore miss every edge with probability $(1-2w)^t\le e^{-2wt}\le1/3$. An edgeless input always passes. The tester uses $2t$ vertex samples and at most $t$ queries and does not require knowing $D$ or the vertex count. The same proof and cost hold for the uniform distribution.
\end{proof}
A tempting general reduction for rational $D$ replaces each vertex by a number of independent clones proportional to its probability and invokes a uniform-distribution tester. This does not automatically preserve the property. For example, the property of being a complete graph is hereditary and extendable. Start with $K_2$ and assign each vertex probability $1/2$; replacing each vertex by two independent clones gives $K_{2,2}$, which is not complete, even though the original input satisfies the property. All sampled clone adjacencies across the two original vertices are correct, but clone pairs corresponding to the same vertex create forbidden missing edges. Thus the proposed black-box transfer fails its perfect-completeness prerequisite.
\subsection{Verification and remaining gaps}
Exact weighted edge masses and miss probabilities were computed on finite graphs, and the clone counterexample was checked. A property-aware treatment of clones could repair this particular example, so it is not a lower bound against all reductions. The edgeless tester gives the same upper-bound order in both models for one property; it neither proves optimality nor compares the sample complexities of all extendable hereditary properties. The source's nonclique forbidden-subgraph question remains untouched.

\section{References and source attribution}
Primary workshop report: \url{https://doi.org/10.4171/owr/2019/19}. Relevant passages read: 2019/19 pp. 1137--1140, VDF access and distance model, qualitative characterization and Problem5.

All retrievals were checked on 11 October 2026. Original question attribution remains in the preserved Sources section. External theorems are attributed and not claimed as new discoveries or fully verified merely by citation.


\paragraph{Integration review.} Codex (GPT-6; exact serving revision not exposed) reviewed the written arguments and expanded layout and source-reference wrappers. Repository metadata and the production web renderer were checked using Python 3.13 and JavaScript with MathJax 3.2.2. This model-assisted review is not independent human proof verification; no PDF compilation or formal proof checking is asserted.

\section{Completion Estimate}
\noindent\textbf{COMPLETION: 0\%}
No substantiated new progress toward the entire original target is claimed. The worked results and failed route are delimited above; a known result, illustrative model, or auxiliary obstruction is not counted as original-target completion.

\section{Final status and attribution}
\textbf{Attempt status: unresolved.} This describes this research attempt, not a new ruling on the source catalogue's status or on external literature claims.
The mathematics and this self-audit were generated by OpenAI \texttt{gpt-6-astra}, reasoning \texttt{ultra}, on 11 October 2026. Codex (GPT-6; exact serving revision not exposed) orchestrated the batch, reviewed the written mathematical arguments, and contributed formatting/repository review. This additional model review is not independent human verification. No human mathematical authorship, independent human proof checking, or formal verification is claimed. The historical preparation paragraphs in the reproduced catalogue body concern that body only.

Contribution prompt A was applied to the complete definition and supplied source record, followed by the adversarial checks of prompt B. Local reproduction files are \texttt{batch022/checks.py}, \texttt{computations.json}, and the per-ID audit JSON. They were not public links at preparation time. Numerical checks are finite and do not prove unrestricted results. No LaTeX engine was found locally; compilation was not performed by this research agent.

\end{document}
